% Einheit 1 von 4 – Wurzelziehen und Lösbarkeit (bank/quadratische-gleichungen/e1.jsonl, zusammenbau v0.1)
%% TODO Zweigzeile Teil 1: was hier gelernt wird (Satz in Schülersprache aus dem Einheitstitel) – Einheit 1
\einheitenkopf[e1]{Einheit 1 von 4 · Wurzelziehen und Lösbarkeit}
\zweigzeile{ab Kl. 9, je nach Buch bis Kl. 10 (am Gymnasium ab Kl. 8) · P10}
\setcounter{aufgabe}{10}

%% TODO Titel: Ich-kann-Satz fehlt in der Bank (Platzhalter: Kettenname) – Nr. 11
%% TODO Anweisung: ein Satz über der ersten Teilaufgabe fehlt in der Bank – Nr. 11
\begin{aufgabe}{Quadratisch oder linear?}
\begin{teile}
\teil $5x - 2 = x^2$ – quadratisch oder linear?\\ \kreuz{quadratisch} \\ \kreuz{linear}
\end{teile}
\end{aufgabe}

%% TODO Titel: Ich-kann-Satz fehlt in der Bank (Platzhalter: Kettenname) – Nr. 12
%% TODO Anweisung: ein Satz über der ersten Teilaufgabe fehlt in der Bank – Nr. 12
\begin{aufgabe}{Wurzelziehen und Lösbarkeit}
\begin{teile}
\teil $(x - 3)^2 = 25$ – wie viele Lösungen?\\ \kreuz{zwei} \\ \kreuz{eine} \\ \kreuz{keine}
\end{teile}
%% TODO Darstellung für schwach fehlt in der Bank (quadratische-gleichungen-e1-k2-s1-v1; Abschnitt „Für schwache Schüler“)
\swz{$\text{x^2 = 81}$}{}
%% TODO Darstellung für schwach fehlt in der Bank (quadratische-gleichungen-e1-k2-s1-v2; Abschnitt „Für schwache Schüler“)
\swz{$\text{x^2 = 25}$}{}
%% TODO Darstellung für schwach fehlt in der Bank (quadratische-gleichungen-e1-k2-s1-v3; Abschnitt „Für schwache Schüler“)
\swz{$\text{x^2 = 100}$}{}
%% TODO Darstellung für schwach fehlt in der Bank (quadratische-gleichungen-e1-k2-s1-v4; Abschnitt „Für schwache Schüler“)
\swz{$\text{x^2 = 169}$}{}
%% TODO Darstellung für schwach fehlt in der Bank (quadratische-gleichungen-e1-k2-s1-v5; Abschnitt „Für schwache Schüler“)
\swz{$\text{x^2 = 196}$}{}
\swfrage{Was bleibt gleich, was ändert sich?}
%% TODO Darstellung für schwach fehlt in der Bank (quadratische-gleichungen-e1-k2-s2-v1; Abschnitt „Für schwache Schüler“)
\swz{Ein quadratisches Beet hat den Flächeninhalt $64\,\text{m}^2$. Seitenlänge?}{}
%% TODO Darstellung für schwach fehlt in der Bank (quadratische-gleichungen-e1-k2-s3-v1; Abschnitt „Für schwache Schüler“)
\swz{$\text{x^2 = 13}$}{}
%% TODO Darstellung für schwach fehlt in der Bank (quadratische-gleichungen-e1-k2-s4-v1; Abschnitt „Für schwache Schüler“)
\swz[3]{$x^2 = -100$ – Lösungen? Begründe.}{}
%% TODO Darstellung für schwach fehlt in der Bank (quadratische-gleichungen-e1-k2-s5-v1; Abschnitt „Für schwache Schüler“)
\swz{$\text{3x^2 = 75}$}{}
%% TODO Darstellung für schwach fehlt in der Bank (quadratische-gleichungen-e1-k2-s6-v1; Abschnitt „Für schwache Schüler“)
\swz[3]{Eine zylinderförmige Regentonne fasst $200$ Liter ($200\,\text{dm}^3$) und ist $8$ dm hoch. Für das Volumen gilt $V = \pi \cdot r^2 \cdot h$ ($r$ Radius, $h$ Höhe). Radius auf zwei Stellen?}{}
%% TODO Darstellung für schwach fehlt in der Bank (quadratische-gleichungen-e1-k2-s7-v1; Abschnitt „Für schwache Schüler“)
\swz{$\text{2x^2 + 5 = 37}$}{}
%% TODO Darstellung für schwach fehlt in der Bank (quadratische-gleichungen-e1-k2-s8-v1; Abschnitt „Für schwache Schüler“)
\swz{$\text{(x - 6)^2 = 16}$}{}
%% TODO Darstellung für schwach fehlt in der Bank (quadratische-gleichungen-e1-k2-s9-v1; Abschnitt „Für schwache Schüler“)
\swz{$\text{(x + 10)^2 = 25}$}{}
%% TODO Darstellung für schwach fehlt in der Bank (quadratische-gleichungen-e1-k2-s10-v1; Abschnitt „Für schwache Schüler“)
\swz{$\text{(x - 9)^2 = 0}$}{}
%% TODO Darstellung für schwach fehlt in der Bank (quadratische-gleichungen-e1-k2-s11-v1; Abschnitt „Für schwache Schüler“)
\swz{Ist $x = -7$ Lösung von $(x + 2)^2 = 25$?}{}
\swz{Parabel $y = (x - 1)^2 - 2$ und Gerade $y = 2$: Wie viele Lösungen hat $(x - 1)^2 - 2 = 2$?}{\begin{ksys}[xmin=-5,xmax=5,ymin=-3,ymax=6,ablesen] \parabel{1}{1}{-2}{} \gerade{0}{2}{} \end{ksys}}
%% TODO Darstellung für schwach fehlt in der Bank (quadratische-gleichungen-e1-k2-s13-v1; Abschnitt „Für schwache Schüler“)
\swz[3]{Gib eine Gleichung der Form $x^2 = c$ ($c$ eine Zahl) an, die keine Lösung hat.}{}
\begin{teile}
\teil Wahr oder falsch? Gib außerdem eine quadratische Gleichung ohne Lösung an. \hfill (P10 2021 OS)\\ A: $(x - 6)^2 = 0$ hat genau eine Lösung. \kreuz{wahr} \kreuz{falsch}\\ B: $7$ und $11$ sind Lösungen von $(x + 9)^2 = 4$. \kreuz{wahr} \kreuz{falsch} Gleichung ohne Lösung: \leerfeld
\end{teile}
\end{aufgabe}

%% TODO Titel: Ich-kann-Satz fehlt in der Bank (Platzhalter: Kettenname) – Nr. 13
%% TODO Anweisung: ein Satz über der ersten Teilaufgabe fehlt in der Bank – Nr. 13
\begin{aufgabe}{Wurzelziehen und Lösbarkeit – Fehler finden, Begründen, Darstellungswechsel, Anwendung}
\swz[3]{Wo ist der Fehler? \rechnung{x^2 &= 225 \\ x &= 15}}{}
\swfrage{Begründe: $x^2 = 10$ hat zwei Lösungen.}
\swa{Zeichne die Gerade $y = 4$ zur Normalparabel $y = x^2$ ein. Lösungen von $x^2 = 4$? x1 ≈ \leerfeld , x2 ≈ \leerfeld}{\begin{ksys}[xmin=-3,xmax=3,ymin=-1,ymax=6] \parabel{1}{0}{0}{} \end{ksys}}
\swz[3]{Ein Stein fällt von einer $80$ m hohen Brücke. Für den Fallweg gilt $s = 5 \cdot t^2$ ($s$ in m, $t$ in s). Nach wie vielen Sekunden schlägt er auf?}{}
\end{aufgabe}

\uebersichtskasten{Wurzelziehen: x² = c hat für positives c zwei Lösungen x₁ = √c und x₂ = −√c, für c = 0 genau eine (x = 0), für negatives c keine – ein Quadrat ist nie negativ. \\ x² = 36 → x₁ = 6, x₂ = −6 \quad x² = 0 → x = 0 \quad x² = −36 → keine Lösung \\ Quadrierte Klammer: (x − d)² = c rückwärts rechnen – erst die Wurzel mit beiden Vorzeichen, dann d hinüberbringen. \\ (x + 4)² = 36 → x + 4 = 6 oder x + 4 = −6 → x₁ = 2, x₂ = −10 \quad (x + 4)² = 0 → x = −4 \\ Probe: Lösung einsetzen, beide Seiten ausrechnen – gleich? (wA)}
